%===================================================================================================
\section{Mechanisms: Expression Capture and Value Extraction} %+++++++++++++++++++++++++++++++++++++
\label{sctn:mechanisms-expression-capture-and-value-extraction} %+++++++++++++++++++++++++++++++++++
%===================================================================================================


% \improve


The output model introduced in the previous section centers on three core
fields: (i) a \emph{tested expression} that communicates what was evaluated in a
harness-independent form --- without needing to know anything about the
surrounding test infrastructure, (ii) an \emph{expected value} that reflects
the specification, and (iii) an \emph{actual value} obtained from executing the
implementation under test. Together, these fields shift the focus of test
feedback from \emph{how} a test is implemented to \emph{what} claim the test
makes about the student program.

\skipM However, standard unit testing APIs typically observe only the
\emph{outcome} of an assertion --- for instance, two already-computed values
passed into an \fsharp|AreEqual| check, or a boolean predicate. The literal
source-level claim that the test author wrote is not preserved. Reconstructing
the \emph{tested expression} field therefore cannot be achieved by merely
forwarding existing assertion parameters. Providing the full test harness would
reintroduce the very problem the model aims to avoid, namely that students must
read and interpret auxiliary test code instead of focusing on the semantic
mismatch.

\skipM A practical observation is that many example-based tests and a large
subset of property-based tests can be normalized to an equality-shaped claim of
the form $\textit{actual} = \textit{expected}$, even if the surrounding test
method constructs inputs, performs intermediate computations, or more complex
tasks. Consequently, the key technical requirement is to capture the \emph{claim
itself} in a form that can later be displayed and analyzed, and then to evaluate
that representation to obtain the expected and actual values for reporting.


%===============================================================================
\subsection{Expression capture via code quotations} %+++++++++++++++++++++++++++
\label{ssctn:expression-capture-via-code-quotations} %++++++++++++++++++++++++++
%===============================================================================


\fs offers code quotations as a structured representation of code that can be
inspected and evaluated later.

\begin{defpanel}
\begin{description}
    \item[{Code quotations:}] a language feature that enables you to generate
    and work with \fs code expressions programmatically. This feature lets you
    generate an abstract syntax tree that represents \fs code. The abstract
    syntax tree can then be traversed and processed according to the needs of
    your application. For example, you can use the tree to generate \fs code or
    generate code in some other language. \cite{CodeQuot}
\end{description}
\end{defpanel}

\noindent A simple example for a code quotation is \fsharp|let expr: Expr<int> =
<@ 1 @>| with \fsharp|1| being the code to be represented, wrapped by the
symbols \texttt{<@} and \texttt{@>}. The code quotation is stored using the
non-generic type \texttt{Expr}. \cite{FsExpr} Consider the following example
expression:

\begin{codepanel}{Example code quotation}
let f = <@ fun (a: int) (b: int) -> a + b @>
\end{codepanel}

%===========================================================
\paragraph{Characterization of represented code.}%
\label{par:characterization-of-represented-code} The code inbetween can be
characterized by a few things: (i) it is an anonymous function or \emph{lambda}
expression, (ii) it is semantically equivalent to \fsharp|fun (a: int) -> (b:
int) -> a + b| and (iii) with that we can also see the explicit type of this
lambda: \fsharp|int -> int -> int|.


%===========================================================
\paragraph{Characterization of expression.}%
\label{par:characterization-the-expression} A look at the string representation
reveals that we are dealing with a nested lambda expression of form:

\begin{center}\fsharp|Lambda (a, Lambda (b, Call (None, op_Addition, [a,
b])))|\end{center}

\noindent which (i) seems to align with the pattern of the original code --- the
types are not printed yet present so type information is not lost, (ii) seems to
be semantically equivalent to the represented code and lastly (iii) type inference
implies the type \fsharp|Expr<int -> int -> int>|.

\skipL In conclusion, the comparison shows that \fsharp|fun (a:int) (b:int) -> a
+ b| denotes a curried function, whereas \fsharp|<@ fun (a:int) (b:int) -> a + b
@>| denotes a quotation of that function expression. Accordingly, the two
expressions are not equivalent at the level of \fs evaluation: the former
evaluates to an executable function value, while the latter evaluates to a
representation of the function's syntax. However, the quotation still captures a
term with the same computational behavior as the unquoted function and since
preserving syntactic structure its crucial advantage is enabling inspection and
transformation of the program as data. At the same time, quotations remain
executable in the sense that they can be evaluated (or partially evaluated)
programmatically and --- where required --- transformed back into executable
function values. Thus, quotations unite syntactic preservation with
computational expressiveness.

\skipM The \fsharp{Microsoft.FSharp.Quotations} namespace offers various active
patterns to analyze expressions, as well as a member \fsharp{Substitute} from
the \fsharp{Expr} type to map variables to new values, which is helpful for
partial evaluations of curried functions.


%===========================================================
\paragraph{Swensen.Unquote library.}
    \label{Swensen.Unquote}
Before going into more detail about the value extraction, it is worth
pointing out that there are already libraries offering testing tools using code
quotations like Swensen.Unquote. \cite{SwenUnq} The following example using
Unquote's \fsharp{test} method illustrates how a quoted equality implicitly
contains all three core fields --- tested expression, expected value, and actual
value:

\begin{outputpanel}{Example \texttt{Swensen.Unquote.test} output}{Example Swensen.Unquote.test}{
Swensen.Unquote.AssertionFailedException: Test failed:\newline\newline
\hlbox[metroBlue]{add 1 2} = \hlbox[metroOrange]{1 + 2}\newline
\hlbox[metroRed]{4} = \hlbox[metroGreen]{3}\newline
false
}
let add a b = a + b + 1
test <@ add 1 2 = 1 + 2 @>
\end{outputpanel}

\noindent In \hlbox[metroBlue]{blue} the tested expression \fsharp|add 1 2|, in
\hlbox[metroRed]{red} the actual value and in \hlbox[metroGreen]{green} the
expected value. Combining code quotations and Swensen.Unquote offers not only
functionality to analyze code via known patterns, but also to easily modify it
via functionalities like: \texttt{eval}, \texttt{reduce}, \texttt{decompile},
\texttt{substitute} or the above mentioned \texttt{test}.
\cite{FsExpr, SwenUnqUserGuide}

While this output makes the mismatch transparent, it also hints at a design
tension: depending on the level of detail, the output may reveal how the
expected value is constructed (here: the right-hand side expression), which can
unintentionally disclose test-oracle information. In \hlbox[metroOrange]{orange}
it is visible that the correct computation seems to be \fsharp|1 + 2|, which
strongly suggest the correct implementation \fsharp|let add a b = a + b|.
Depending on the usage, information may leak that is not intended for students.

\skipM While offering a good foundation to suit the usecase some optimization
are needed which I will mention later.


%===============================================================================
\subsection{Value extraction for unit tests and property based tests} %+++++++++
\label{ssctn:value-extraction-for-unit-tests-and-property-based-tests} %++++++++
%===============================================================================


While quotations preserve the structure of a test claim, Assertify still needs
to \emph{separate} the claim into the parts that correspond to the output model
(i.e., tested expression, expected value, and actual value) and then
\emph{evaluate} those parts. This is achieved by restricting the supported
shape of claims to a normal form that is both common in introductory exercises
and directly compatible with the reporting scheme: an equality-shaped claim of
the form $\textit{actual} = \textit{expected}$.

\skipM In practice, such claims are represented as quoted boolean expressions.
At runtime, Assertify and Checkify inspect these quotations using standard
quotation patterns. In the intended case, the quotation contains a call to the
equality operator, which can be deconstructed into its left-hand side (the
implementation/result under test) and right-hand side (the reference
implementation or result). This deconstruction step forms the bridge from quoted
code to the three-field output model: the \emph{tested expression} is rendered
from the left-hand side (with concrete inputs substituted, if necessary), while
\emph{actual} and \emph{expected} are obtained by evaluating the left and right
subexpressions respectively.


%===========================================================
\paragraph{Unit tests.}
\label{par:value-extraction-unit-tests}
For ordinary unit tests, the quoted claim already contains concrete values.
Accordingly, extraction proceeds directly by deconstructing the equality-shaped
quotation into its two subexpressions and evaluating them independently.
Conceptually, this maps
\begin{align*}
\texttt{actualExpr} &= \texttt{\hlbox[metroBlue]{add 1 2}} &
\texttt{expectedExpr} &= \texttt{\hlbox[metroOrange]{1 + 2}}\\
\llbracket\,\texttt{actualExpr}\,\rrbracket &= \texttt{\hlbox[metroRed]{4}} &
\llbracket\,\texttt{expectedExpr}\,\rrbracket &= \texttt{\hlbox[metroGreen]{3}}
\end{align*}
where $\llbracket\,\cdot\,\rrbracket$ denotes evaluation of a quotation.

\skipM This also ensures that no solution-relevant information is leaked to
students in raw syntactic form. For example, in a claim such as
\fsharp|add 1 2 = 1 + 2|, the critical right-hand side is not exposed as source
code in the feedback, but only through its evaluated result. This supports the
didactic goal of revealing enough information to diagnose the error while
avoiding direct disclosure of the reference solution.

\skipM This also assures that no information is leaked to students,
because the critical right hand side \fsharp|1 + 2| in \ref{code:Example
Swensen.Unquote.test} would only be visible to the student in evaluated form.

\skipM With that a stable basis for reporting is provided, because it avoids
reliance on ad-hoc parsing of framework error strings and instead operates on
the structured representation of the asserted claim. If the quotation does not
match the expected equality pattern, the mechanism cannot reliably derive
\emph{expected} and \emph{actual} values at this point of implementation. In
that case, Assertify falls back to reporting the captured expression and a
generic failure message. This limitation is deliberate: extending support to
arbitrary assertion forms is possible, but requires separate, assertion-specific
deconstruction rules (e.g., inequalities, exception checks, boolean predicates).


\begin{algorithm}
\caption{Assertify.Test}
\begin{algorithmic}[0]
    \Function{Test}{\texttt{expr: Expr<bool>, ?message: string}}
        \If{\texttt{message} $=$ \texttt{null}}
            \State\texttt{message} $\gets$ \texttt{defaultMsg}
        \EndIf
        \Try
            \If{$\llbracket$ expr $\rrbracket$}
                \State\Return\texttt{success}
            \Else
                \State\Return\Call{ConstructErrorMsg}{\texttt{expr}, \texttt{message}}
            \EndIf
        \With
            \State\Return\Call{DefaultErrorMsg}{\texttt{expr}, \texttt{message}}
            % \Comment{\texttt{Failed to compute $\llbracket$ expr $\rrbracket$ or}}
            % \Comment{\texttt{to extract a left/right hand side}}
    \EndFunction
\end{algorithmic}
\end{algorithm}


\begin{algorithm}
\caption{Checkify.Check}
\begin{algorithmic}[0]
    \Function{Check}{\texttt{expr: Expr, check: Expr -> 'Testable}}
        \State \texttt{capture runner}
        \Try
            \State\texttt{FsCheck.Check (check expr)}
            \Comment{\texttt{Successful test}}
            \State\Return\texttt{success}
        \With{Hallo}
        \Try{\texttt{add 1 2}}
            \State\texttt{Test 1 2 3}
    \EndFunction
\end{algorithmic}
\end{algorithm}


%===================================================================================================
%TODO TODO TODO TODO TODO TODO TODO TODO TODO TODO TODO TODO TODO TODO TODO TODO TODO TODO TODO TODO
%===================================================================================================


% Concrete implementation details?

\vfill
\pagebreak

\improve
\pagebreak


%
\[
\textit{actualValue} = \lbrack \textit{actualExpr} \rbrack
\qquad\text{and}\qquad
\textit{expectedValue} = \lbrack \textit{expectedExpr} \rbrack
\]
%
where $\lbrack \cdot \rbrack$ denotes evaluation of a quotation. In the
implementation, evaluation is performed using the quotation evaluation facilities
provided by the \fs runtime and by Swensen.Unquote.

\skipM A practical complication is that evaluation may itself fail. This happens,
for example, when the student implementation throws an exception, when a branch
contains a remaining placeholder such as \fsharp|failwith "TODO"|, or when a
quoted expression cannot be evaluated due to missing runtime context. For novice
feedback, a raw evaluation exception would reintroduce the problem of
unstructured output. Assertify therefore treats evaluation failure as a first-class
case and produces a controlled placeholder (e.g., ``non-evaluatable'') rather than
propagating the exception as the primary message. The full technical details of
these fallbacks are discussed later with the concrete helper functions.

%===============================================================================
\subsubsection{Example-based tests: closed expressions} %++++++++++++++++++++++++
\label{sec:value-extraction-example}
%===============================================================================

In example-based unit tests, the quoted claim is typically \emph{closed}: all
inputs are written explicitly in the test case (e.g., \fsharp|digitSum 123N = 6N|).
This means that value extraction can proceed immediately after capture:
deconstruct the equality, evaluate both sides, and render the tested expression.

\skipM For reporting, Assertify distinguishes two closely related representations:
(i) a direct rendering of the left-hand side as the tested expression (to answer
``what was executed?'') and (ii) the evaluated values of both sides (to answer
``what was expected and what was produced?''). In this setting, no additional
context is required beyond the captured quotation itself.

%===============================================================================
\subsubsection{Property-based tests: open expressions and captured inputs} %++++++
\label{sec:value-extraction-property}
%===============================================================================

Property-based tests differ in one crucial aspect: the claim is generally
\emph{open} because it depends on generated inputs. Instead of quoting a closed
boolean, a property-based test typically captures a quoted lambda expression,
such as \fsharp|<@ fun n -> sortedDigits n = sortedDigitsSolution n @>|. The
quotation therefore describes the \emph{shape} of the claim, but evaluation and
rendering require concrete input values.

\skipM FsCheck provides these values only when a counterexample is found. In
addition, FsCheck may \emph{shrink} a failing input to a simpler one. For the
purposes of novice-oriented feedback, the most relevant inputs are therefore the
final counterexample values after shrinking, because they are the smallest (and
usually easiest to reason about) values that still falsify the property.

%===============================================================================
\subsubsection{Runner-assisted extraction with \texttt{CaptureRunner}} %+++++++++
\label{sec:runner-assisted-extraction}
%===============================================================================

To make counterexample inputs available to the reporting layer, Checkify wraps
the configured FsCheck runner with a custom runner (\texttt{CaptureRunner}). The
runner intercepts the test result at the end of execution. When FsCheck reports a
failure, the runner stores both the original failing arguments and the final
shrunk arguments.

\skipM The stored inputs are then used in two places:

\begin{enumerate}
    \item \textbf{Instantiation for rendering.} The open quotation (a lambda) is
    instantiated by substituting the captured counterexample values for its
    parameters. This yields a quotation that corresponds to the concrete failing
    call, which can be rendered as the \emph{tested expression} shown to the
    student.

    \item \textbf{Instantiation for evaluation.} The same substitution step
    yields a concrete equality claim whose left and right subexpressions can be
    deconstructed and evaluated to obtain \emph{actual} and \emph{expected}.
\end{enumerate}

\skipM In Checkify, this instantiation is implemented as iterative application of
arguments to a quoted lambda, using quotation substitution. The application
function traverses nested lambdas and replaces parameter occurrences with
\fsharp|Expr.Value| nodes constructed from the captured objects. This keeps the
mechanism uniform across different arities and preserves the key invariant:
\emph{the displayed tested expression corresponds to the failing counterexample
that FsCheck reports}.

%===============================================================================
\subsubsection{Reduction traces: diagnostic benefit vs.\ oracle leakage} %++++++++
\label{sec:reduction-traces}
%===============================================================================

Some quotation-based tools provide additional context by showing a reduction
trace, i.e., intermediate simplifications of an expression until a final form is
reached. This can make the connection between a complex claim and its final
mismatch more explicit, particularly for novices who struggle to mentally execute
expressions.

\skipM At the same time, reduction traces can reveal how an expected value was
constructed (for instance by exposing the full right-hand side computation),
which may disclose test-oracle information. Since Assertify targets an
exercise-oriented setting, it treats such detail as optional and subject to
verbosity control rather than as a mandatory default. This makes it possible to
trade diagnostic transparency against pedagogical constraints depending on the
context in which the feedback is used.

\begin{lstlisting}[caption={Example of a reduction trace (adapted from the Unquote user guide)},label={lst:unquote-trace}]
Test failed:

(1 + 2) / 3 = 2
3 / 3 = 2
1 = 2
false
\end{lstlisting}

%===============================================================================
\subsubsection{Summary} %+++++++++++++++++++++++++++++++++++++++++++++++++++++++
\label{sec:mechanisms-summary}
%===============================================================================

Expression capture and value extraction are the technical prerequisites for
producing structured test feedback in the form of \emph{tested expression},
\emph{expected value}, and \emph{actual value}. Quotations enable the capture of
asserted claims as structured data; quotation patterns allow the system to
identify and deconstruct equality-shaped claims; and evaluation of the resulting
subexpressions yields the runtime values required for reporting.

\skipM For property-based testing, quotations alone are not sufficient because
claims are open with respect to generated inputs. Checkify therefore integrates
at the runner level: \texttt{CaptureRunner} records the original and shrunk
counterexample arguments so that the captured quotation can be instantiated,
rendered, and evaluated in the concrete context of the failing assignment. The
implementation details of these helper functions and their formatting choices are
covered later in this chapter.


%===================================================================================================
% EOF ++++++++++++++++++++++++++++++++++++++++++++++++++++++++++++++++++++++++++++++++++++++++++++++
%===================================================================================================